% Part: turing-machines
% Chapter: machines-computations
% Section: well-behaved-machines

\documentclass[../../../include/open-logic-section]{subfiles}

\begin{document}

\olfileid{tur}{mac}{dis}
\olsection{शिस्तबद्ध यंत्रे}

\begin{explain}
\olref[hal]{sec} मध्ये आपण एकच निर्दिष्ट थांबण्याची अवस्था~$h$
असलेली ट्यूरिंग यंत्रे पाहिली: अशी यंत्रे थांबलीच तर~$h$
अवस्थेतच थांबतात. म्हणून एका थांबण्याच्या अवस्थेची यंत्रे, आपण
सर्वसाधारण ट्यूरिंग यंत्रांना जितकी मुभा देतो त्यापेक्षा अधिक
``शिस्तबद्ध'' आहेत. ट्यूरिंग यंत्रांच्या वर्तनावर आणखीही निर्बंध
घालता येतील. उदाहरणार्थ, घर~$0$ वरील फितीच्या डाव्या टोकाचे
चिन्ह पुसण्यास किंवा घर~$0$ वरून डावीकडे जाण्याचा प्रयत्न
करण्यास आपण त्यांना मनाई केलेली नाही. (आपल्या व्याख्येनुसार
अशा वेळी शीर्ष घर~$0$ वरच राहते; इतर व्याख्यांनुसार यंत्र
थांबते.) तसेच यंत्र थांबलेच तर घर~$1$ वर थांबेल, असे गृहीत
धरता आले तर काही वेळा सोयीचे ठरते.
\end{explain}

\begin{defn}\ollabel{defn:disciplined}
ट्यूरिंग यंत्र~$M$ \emph{शिस्तबद्ध} असते तेव्हा आणि केवळ तेव्हाच,
जेव्हा
\begin{enumerate}
    \item त्याला एकच निर्दिष्ट थांबण्याची अवस्था~$h$ असते,
    \item ते थांबलेच तर घर~$1$ वाचत असताना थांबते,
    \item ते घर~$0$ वरील $\TMendtape$ चिन्ह कधीही पुसत नाही, आणि
    \item ते घर~$0$ वरून डावीकडे जाण्याचा कधीही प्रयत्न करत नाही.
\end{enumerate}
\end{defn}

\begin{explain}
कोणतेही ट्यूरिंग यंत्र त्याच वर्तनाचे पण निर्दिष्ट थांबण्याची
अवस्था असलेले यंत्र बनवता येते, हे आपण आधी पाहिले. यासाठी
नवी अवस्था~$h$ जोडायची आणि मूळ~$\delta$ ज्या
$\tuple{q,\sigma}$ जोडीवर अपरिभाषित आहे, त्या प्रत्येक जोडीसाठी
$\delta(q, \sigma) = \tuple{h, \sigma, N}$ ही सूचना जोडायची.
पुरावा किचकट असला तरी कोणतेही ट्यूरिंग यंत्र~$M$ अशा
शिस्तबद्ध ट्यूरिंग यंत्रात~$M'$ बदलता येते की जे त्याच आदानांवर
थांबते आणि तेच प्रदान देते. उदाहरणार्थ, यंत्र थांबत असताना
घर~$1$ वर नसेल, तर शीर्ष फितीच्या डाव्या टोकाचे चिन्ह सापडेपर्यंत
डावीकडे नेणाऱ्या, त्यानंतर एक घर उजवीकडे नेणाऱ्या आणि मग यंत्र
थांबवणाऱ्या काही सूचना जोडता येतात. इतर अटी कशा पूर्ण करता
येतील याचा विचार तुम्ही करा.
मूळ यंत्राने डाव्या टोकाचे चिन्ह पुसले किंवा तेच चिन्ह दुसऱ्या घरावर
लिहिले असेल, तर एवढी साधी शोध-पद्धत पुरेशी नाही; सर्वसाधारण
रूपांतराच्या पुराव्यात हा प्रसंग वेगळा हाताळावा लागतो.
%READERNOTE{OLTUR-006 — गोठवलेल्या मजकुरात शीर्ष डाव्या टोकाचे चिन्ह सापडेपर्यंत डावीकडे नेण्याची पद्धत दिली आहे. पण आधीची यंत्र-व्याख्या ते चिन्ह पुसण्यास किंवा इतर घरांत लिहिण्यास मनाई करत नाही. त्यामुळे हा केवळ मर्यादित उदाहरणार्थ उपाय आहे; सर्वसाधारण शिस्तबद्ध रूपांतरासाठी अधिक अचूक चिन्ह-व्यवस्था आवश्यक आहे.}
\end{explain}

\begin{ex}
    \olref{fig:adder-disc} मध्ये आपण \olref[una]{ex:adder}
    मधील बेरीज-यंत्राचे शिस्तबद्ध यंत्रात रूपांतर करतो.
    \begin{figure}\[
\begin{tikzpicture}[->,>=stealth',shorten >=1pt,auto,node distance=2.8cm,
                    semithick]
  \tikzstyle{every state}=[fill=none,draw=black,text=black]

  \node[initial,state]         (A)              {$q_0$};
  \node[state]         (B) [right of=A] {$q_1$};
  \node[state]         (C) [below right of=B] {$q_2$};
  \node[state]         (D) [below left of=C] {$q_3$};
  \node[state]         (H) [left of=D] {$h$};

  \path (A) edge node {\TMtrans{\TMblank}{\TMstroke}{\TMstay}} (B)
           edge [loop below] node {\TMtrans{\TMstroke}{\TMstroke}{\TMright}} (B)
        (B) edge [loop below] node {\TMtrans{\TMstroke}{\TMstroke}{\TMright}} (B)
            edge node {\TMtrans{\TMblank}{\TMblank}{\TMleft}} (C)
        (C) edge node {\TMtrans{\TMstroke}{\TMblank}{\TMleft}} (D)
        (D) edge [loop above] node {\TMtrans{\TMstroke}{\TMstroke}{\TMleft}} (D)
        edge node {\TMtrans{\TMendtape}{\TMendtape}{\TMright}} (H);
\end{tikzpicture}
\]\caption{शिस्तबद्ध बेरीज-यंत्र}
\ollabel{fig:adder-disc}
\end{figure}
\end{ex}

\begin{prop}\ollabel{prop:disciplined} प्रत्येक ट्यूरिंग यंत्र~$M$
साठी असे शिस्तबद्ध ट्यूरिंग यंत्र~$M'$ असते की $M$~प्रदान~$O$
देऊन थांबत असेल तर ते देखील प्रदान~$O$ देऊन थांबते; आणि
$M$~थांबत नसेल तर ते देखील थांबत नाही. विशेषतः, ट्यूरिंग
यंत्राला संगणित करता येणारे कोणतेही $f\colon\Nat^n \to \Nat$
फलन शिस्तबद्ध ट्यूरिंग यंत्रालाही संगणित करता येते.
\end{prop}

\begin{prob}\label{tur:mac:dis:prob:disc-succ}
$f(x) = x+1$ संगणित करणारे शिस्तबद्ध यंत्र द्या.
\end{prob}


\begin{prob}\label{tur:mac:dis:prob:copier}
$\TMstroke^n$ आदानावर सुरू केल्यावर
$\TMstroke^n \concat \TMblank \concat \TMstroke^n$ हे प्रदान
देणारे शिस्तबद्ध यंत्र शोधा.
\end{prob}

\end{document}
